Une personne dévêtue, dont la peau (émissivité \num{0.97} dans l'infrarouge) est à \qty{\v{Ts}}{\degreeCelsius} et dont la surface corporelle est de \qty{\v{S}}{m^2}, se trouve dans une pièce dont les murs sont à \qty{\v{Te}}{\degreeCelsius}.

\begin{enumerate}[label=\alph*)]
    \item Déterminez la puissance nette que la personne perd par rayonnement.
    \item Quelle énergie perd-elle ainsi en \qty{\v{t}}{h} ? Exprimez-la en kilocalories alimentaires ($\qty{1}{kcal} = \qty{4186}{J}$).
    \item Pourquoi a-t-on froid près d'une grande fenêtre en hiver, même si l'air de la pièce est chaud ?
\end{enumerate}
